% Template for post or presentation abstracts submitted to ILAS 2022
% 1. Fill in the presenter's information (name, affiliation, email
% address) and talk information (title of presentation, abstract).
% 2. Compile to make sure there are no errors.
% 3. Save the .tex file for your abstract with a distinctive name,
% ideally "FamilynameGivenname.tex".
% 4. Upload your .tex file to https://clr.ie/129557
%
% * Please keep your LaTeX commands simple, and avoid the use of
%    user-defined macros.
% * Include details of co-authors and funding acknowledgements after the abstract.
% * Upload only the LaTeX file (with .tex extension).
% * Questions: Contact Niall Madden (Niall.Madden@NUIGalway.ie)

\documentclass[a4paper]{amsart}
\usepackage{amssymb}
\usepackage{hyperref}
\pagestyle{empty}
\date{}

%% IGNORE THE NEXT 4 LINES
\makeatletter %
\def\labelenumi{\theenumi.}
\def\theenumi{\@arabic\c@enumi}
\makeatother
%% STOP IGNORING NOW

\begin{document}

\title{On the algebraic structure \\ of the copositive cone}
\author{Roland Hildebrand} %Presenting author only
\address{Univ. Grenoble Alpes, CNRS, Grenoble INP, LJK, 38000 Grenoble, France} % Affiliation of presenting author
\email{roland.hildebrand@univ-grenoble-alpes.fr} % Email address of presenting author only

\maketitle

% The abstract  ***no more than one page***

A closed convex cone can be decomposed into a disjoint union of interiors of its faces. This well-known facial decomposition yields a lot of information on the structure of the cone. However, in general there are infinitely many faces, and for some purposes this decomposition is too fine. Some cones admit a coarser, finite decomposition which unites faces which are of the same type. For example, the cone of positive semi-definite matrices of size $n$ decomposes into $n+1$ relatively open manifolds, each of which contains positive semi-definite matrices of constant rank and which are themselves unions of interiors of similar faces.

We propose such a finite decomposition for the copositive cone ${\mathcal COP}^n$. The components of the decomposition are parameterized by the \emph{extended minimal zero support set}. This means that each component $S_{\mathcal E}$ is composed of copositive matrices $A$ with the same extended minimal zero support set ${\mathcal E}$. This set is a collection of pairs ${\mathcal E} = (I_{\alpha},J_{\alpha})_{\alpha = 1,\dots,|{\mathcal E}|}$, where $\alpha$ enumerates the minimal zeros $u_{\alpha}$ of $A$, $I_{\alpha}$ is the support of the minimal zero $u_{\alpha}$, and the index set $J_{\alpha} \supset I_{\alpha}$ consists of those indices $j \in \{1,\dots,n\}$ such that $(Au_{\alpha})_j = 0$.

The set $S_{\mathcal E}$ lies in a real-algebraic variety $Z_{\mathcal E}$ which is given by a finite number of polynomial equalities, namely those equivalent to the rank-deficiency of the sub-matrix $A_{I_{\alpha} \times J_{\alpha}}$. Our main result states that for every $A \in {\mathcal COP}^n$ with extended minimal zero support set ${\mathcal E}$, there exists a neighbourhood $U$ of $A$ in the space of real symmetric matrices such that $U \cap Z_{\mathcal E} \subset S_{\mathcal E}$, i.e., $S_{\mathcal E}$ is open in $Z_{\mathcal E}$. Thus the polynomial equalities cited above fully determine the local structure of $S_{\mathcal E}$.

\begin{thebibliography}{1}
\bibitem{Hildebrand20a}
Roland Hildebrand.
\newblock On the algebraic structure of the copositive cone.
\newblock {\em  Optim. Lett.}  14(8):2007–2019,  (2020).
\end{thebibliography}


\end{document}


% Please leave the next bit alone
%%% Local Variables:
%%% mode: latex
%%% TeX-master: "../ILAS-2022"
%%% End:
